\subsection{Hilbert Basis Theorem}

\begin{theorem} \label{thm:hilbert-basis-theorem} \hlt{Hilbert Basis Theorem}\\
Let $A$ be commutative with identity. If $A$ is Noetherian, then $A[x]$ is Noetherian.
\begin{proof}
It suffices to show every ideal $\mathcal{J}\lhd A[x]$ is f.g. If $\mathcal{J}=0$, done.\\
For $g\neq0$, write $\operatorname{lc}(g)$ and $\operatorname{lt}(g)$. Set $L=\langle\operatorname{lc}(g):g\in\mathcal{J}, g\neq0\rangle\lhd A$. Since $A$ is Noetherian, choose $a_1,\ldots,a_m$ generating $L$ and $f_i\in\mathcal{J}$ with $\operatorname{lc}(f_i)=a_i$. Replacing $f_i$ by powers of $x$ times $f_i$, assume $\deg f_i=d$ for all $i$.\\
For $0\leq k<d$, set $L_k=\langle\operatorname{lc}(g):g\in\mathcal{J}, \deg g=k\rangle$ and choose finite $P_k\subseteq\mathcal{J}$ with $\deg p=k$ for $p\in P_k$ and $\{\operatorname{lc}(p):p\in P_k\}$ generating $L_k$. Put $G=\{f_1,\ldots,f_m\}\cup\bigcup_{k=0}^{d-1}P_k$ and $\mathcal{A}=\langle G\rangle\lhd A[x]$. Then $\mathcal{A}\subseteq\mathcal{J}$.\\
For $g\in\mathcal{J}\setminus\mathcal{A}$ of minimal degree, set $r=\deg g$. If $r<d$, write $\operatorname{lc}(g)=\sum_{p\in P_r}c_p\operatorname{lc}(p)$ and set $h=\sum_{p\in P_r}c_pp$.\\
If $r\geq d$, write $\operatorname{lc}(g)=\sum_{i=1}^{m}c_ia_i$ and set $h=\sum_{i=1}^{m}c_ix^{r-d}f_i$. In either case $h\in\mathcal{A}$ and $\operatorname{lt}(h)=\operatorname{lt}(g)$.\\
Then $g-h\in\mathcal{J}\setminus\mathcal{A}$ and $\deg(g-h)<\deg g$, contradiction. Thus $\mathcal{J}=\mathcal{A}$ is f.g.\\
\hlt{Structural Commentary}: Ideals in $A[x]$ are controlled by finitely many leading coefficients in $A$.\\
\hlt{Key Insight}: Match the leading term, subtract, and induct on degree.
\end{proof}
\end{theorem}
